\section{Multi-Modal Applications}
\label{sec:mutil-modal}

The integration of multi-modal data is indispensable for advancing the capabilities and applications of multi-agent with LLMs.
\ours is designed to seamlessly support various data modalities, leveraging the diverse inputs and outputs that contemporary LLMs can process and produce.

\begin{figure}[ht]
    \centering
    \includegraphics[width=0.7\linewidth]{figures/modals_2.png}
    \caption{The generation, storage, and transmission of Multi-modal data in \ours.}
    \label{fig:modals}
\end{figure}

\paragraph{Management of Multi-Modal Data}
In a running \ours application, the lifecycle of multi-modal data is carefully managed. This management includes the generation, transmission, and storage of multi-modal data—all facilitated through a decoupled architecture using URLs and a local file manager system. 
\figref{fig:modals} exemplifies this process, including data originating from user inputs or model generations, data storage and retrieval, and data sharing.

\begin{itemize}
    \item 
    \textbf{Multi-modal data generation}:
    There are two primary sources of multi-modal data in \ours.
    One source is simply the locally stored multi-modal files, which can be used by either user proxy agents or general agents with access to the local file system.
    Another source is the model-modal content generation models.
    Our model APIs and the model wrappers integrate the most popular multi-modal models, such as the text-to-image content generation models like OpenAI's DALL-E, and conversely, the image-to-text image analysis models, e.g., GPT-4V.
    Besides the built-in APIs, developers can introduce their favorite multi-modal models and customize their own model wrappers, with our ready-to-use examples as the starting points.
    This customization process is streamlined in \ours and benefits from our modular design, allowing developers to connect their multi-modal services with minimal effort.
    \item 
    \textbf{Multi-modal data storage}:
    As mentioned above, multi-modal data in the multi-agent application can be either from ready-to-use local files or generated by multi-modal models.
    When a multi-modal model wrapper is invoked to generate multi-modal data, it first saves the data locally with the help of the file manager and returns a local URL when it receives multi-modal data from the model API service. 
    \item
    \textbf{Multi-modal data transmission}:
    \ours simplifies the process of multi-modal data sharing between agents by allowing agents to encapsulate local or remote URLs in multi-modal messages to indicate the actual storage locations of the data.
    The receiver agents can load the multi-modal data through the URLs when ready to process those.
\end{itemize}

The benefits of introducing URLs in the messages when agents share multi-modal data are three-fold.
Firstly, it can minimize the message size to avoid potential errors or delays because of the network bandwidth and enable the receiver agent to load the data on demand.
Secondly, if there is other text information in the message, the downstream agents can potentially prioritize or parallel the processing of the text information to/and the processing of multi-modal information.
Last but not least, such URL-attached messages can also facilitate the multi-modal data demonstration, which will be introduced in the following section.

\paragraph{Multi-Modal Interaction Modes}
With the implementation of URL-attached messages, \ours empowers users to interact with multi-modal systems via accessible interfaces such as terminal and web UI.  \figref{fig:modal_webui} showcases the user's ability to interact with multi-modal data within interaction modes. 
In the terminal, users can conveniently access locally stored data by activating the provided URLs. The web UI further enhances user experience by providing an intuitive platform to view and analyze multi-modal content, aligning with the expectations of modern web applications.

\paragraph{} Through \ours, developers are equipped to tailor model API services and wrappers to their individual needs, forge applications that handle diverse data modalities, and provide users with the necessary tools to engage with multi-modal agents effectively. This comprehensive support for multi-modal applications positions \ours as a versatile and powerful framework for harnessing the full potential of multi-agent LLMs, broadening the horizons for developers and researchers alike in creating sophisticated and interactive AI systems.
